% Exemplo: fractais com a biblioteca lindenmayersystems do TikZ: o floco de neve de Koch e uma
% planta, descritos por regras de reescrita (sistemas de Lindenmayer).
% Compilar: latexmk -pdf fractal.tex              Testado com: TikZ 3.1, TeX Live 2024 ou mais recente
%%BEGIN:fractalpreambulo
\documentclass[12pt, a4paper]{article}
\usepackage[T1]{fontenc}
\usepackage[brazilian]{babel}
\usepackage{tikz}
\usetikzlibrary{lindenmayersystems}
%%END:fractalpreambulo

\begin{document}
\pagestyle{empty}

%%BEGIN:fractalkoch
\pgfdeclarelindenmayersystem{Koch}{%
  \rule{F -> F-F++F-F}}              % cada segmento vira quatro
\begin{tikzpicture}
  \draw[blue!70!black, l-system={Koch, axiom=F++F++F, order=4,
        step=0.8mm, angle=60}]
    lindenmayer system -- cycle;
\end{tikzpicture}
%%END:fractalkoch
\hfill
%%BEGIN:fractalplanta
\pgfdeclarelindenmayersystem{Planta}{%
  \rule{X -> F[+X]F[-X]+X}           % [ e ] guardam e recuperam a posição
  \rule{F -> FF}}
\begin{tikzpicture}
  \draw[green!50!black, rotate=90, l-system={Planta, axiom=X, order=5,
        step=1.1mm, angle=25}]
    lindenmayer system;
\end{tikzpicture}
%%END:fractalplanta

\end{document}
